\documentclass[11pt]{article}
\usepackage{manual_docs_style}
\externaldocument{build/environment_profiles}[environment_profiles.pdf]
\externaldocument{build/candidate_generation}[candidate_generation.pdf]
\externaldocument{build/simulation_stage}[simulation_stage.pdf]

\begin{document}

\docTitle{Stage: Cheap Filter Metrics}
\phantomsection\label{docstart:cheap_filter_metrics_stage}

This stage computes low-cost detectability, numerical-validity, and SNR metrics from materialized simulation runs.
It rejects failed, non-detectable, numerically invalid, and low-SNR candidates.
It does not prove distinguishability from all alternative root-cause classes and does not run optimized counterfactual search.
Final benchmark selection is performed later by combining \code{cheap_filter_pass} with \code{surrogate_filter_pass}.

\section*{Inputs}

For each \termSimulator, the stage reads:

\begin{itemize}
\item \code{plans/<policy_id>/run_nodes.jsonl}: canonical run recipes and \code{run_id} values.
\item \code{plans/<policy_id>/run_edges.jsonl}: baseline/intervention/time-zero relationships.
\item \code{runs/model_record.json}: runtime status and recipe hashes.
\item \code{runs/<run_id>/data.parquet}: materialized trajectories.
\item \code{experiment_config.json}: observed signals, signal types, and detectability thresholds.
\item \code{noise_adder.py}: \termSimulator-specific noise profiles and SNR analysis.
\item \code{detectability_specific_environment.py}: optional environment-specific detectability veto hook.
\end{itemize}

During migration, the stage may also read legacy \code{model_run_specs.json}, but new code should use the resolved run graph.

\section*{Run Relationships}

For each intervention node, the stage uses \code{run_edges.jsonl} to find:

\begin{itemize}
\item the matched no-intervention baseline run;
\item the matched time-zero/static-change probe run;
\item the changed root-cause parameter;
\item the intervention time and post-intervention value.
\end{itemize}

These relationships are not inferred from filename or identifier structure.

\section*{Cheap-Filter Computation}

For each successful intervention run, compute detectability against:

\begin{itemize}
\item its matched baseline run;
\item its matched time-zero/static-change run.
\end{itemize}

Both comparisons load the dataframes, assert aligned timestamps, apply \termSimulator-specific noise when requested, preprocess channels according to the environment and paper appendix, compute SRD threshold crossings, compute RMS distances, and compute effect-to-noise SNR values.

The generic output includes:

\begin{itemize}
\item \code{mean_rms_distance_clean_dirty}
\item \code{mean_rms_distance_clean_baseline}
\item \code{mean_SNR}
\item \code{first_diff}
\item \code{detectable}: \code{"yes"}, \code{"no"}, or \code{"error"}
\end{itemize}

For the comparison against the no-intervention baseline, the stage may additionally call:

\begin{DocCode}
detectability_specific_environment.py::is_detectable(...)
\end{DocCode}

If this hook returns \code{"no"}, the intervention fails the cheap filter even if the generic RMS/SRD checks pass.

\section*{Pass Rules}

An intervention has \code{cheap_filter_pass = true} only if:

\begin{itemize}
\item it has status \code{success} in \code{model_record.json};
\item it is numerically valid and has all required observable signals;
\item its comparison against the matched baseline has \code{detectable == "yes"};
\item its comparison against the matched time-zero/static-change run has \code{detectable == "yes"};
\item the relevant noise levels satisfy configured SNR thresholds;
\item no \termSimulator-specific low-cost veto hook rejects it.
\end{itemize}

A family has \code{family_cheap_filter_pass = true} when it has cheap-filter-passing children for at least the configured minimum number of distinct root-cause parameters.
The default minimum is three distinct parameters unless a policy or selection stage specifies a stricter rule.

\section*{Output}

The preferred output location is:

\begin{DocCode}
models/simulink/<simulator_id>/metrics/<policy_id>/cheap_filter_metrics.jsonl
\end{DocCode}

During migration, the same payload may also be written to \code{eligibility_metrics.jsonl} as a backward-compatible alias.
New documentation and new code should use \code{cheap_filter_metrics.jsonl}.

A legacy aggregate file may also be written during migration:

\begin{DocCode}
models/simulink/<simulator_id>/runs/eligibility_metrics.json
\end{DocCode}

\section*{Preferred JSON Lines Shape}

Each line of \code{cheap_filter_metrics.jsonl} records one run or one family summary.

Example intervention record:

\begin{DocCode}
{
  "record_type": "run",
  "policy_id": "production_v1",
  "family_id": "fam_de4d7c0a5335e9005510cc2442d0971e",
  "run_id": "run_7fd9daa0055fb3f412ef5e0ab8d4c5ba",
  "kind": "intervention",
  "ground_truth_label": "gravity",
  "baseline_run_id": "run_baseline_...",
  "time0_run_id": "run_time0_...",
  "changed_parameter": "gravity",
  "cheap_filter_pass": true,
  "detectability": {
    "vs_baseline": {"detectable": "yes"},
    "vs_time0_baseline": {"detectable": "yes"}
  },
  "snr": {
    "none": null,
    "low": {"mean_snr_db": 18.2},
    "high": {"mean_snr_db": 9.7}
  },
  "failure_reasons": []
}
\end{DocCode}

Example baseline or no-intervention record:

\begin{DocCode}
{
  "record_type": "run",
  "policy_id": "production_v1",
  "family_id": "fam_de4d7c0a5335e9005510cc2442d0971e",
  "run_id": "run_baseline_...",
  "kind": "baseline",
  "ground_truth_label": "no_intervention",
  "cheap_filter_pass": true,
  "detectability": null,
  "snr": {
    "none": null,
    "low": {"mean_snr_db": 22.4},
    "high": {"mean_snr_db": 12.1}
  },
  "failure_reasons": []
}
\end{DocCode}

Example family record:

\begin{DocCode}
{
  "record_type": "family",
  "policy_id": "production_v1",
  "family_id": "fam_de4d7c0a5335e9005510cc2442d0971e",
  "baseline_run_id": "run_baseline_...",
  "family_cheap_filter_pass": true,
  "cheap_filter_pass_parameters": ["gravity", "mass", "drag_coeff"]
}
\end{DocCode}

\section*{Validation Rules}

\begin{itemize}
\item Every metric record must refer to a \code{run_id} present in \code{run_nodes.jsonl}.
\item Every intervention metric must have a matched baseline and time-zero/static-change run from \code{run_edges.jsonl}.
\item Only successful materialized runs may have \code{cheap_filter_pass = true}.
\item The cheap-filter stage must not change run recipes or recipe hashes.
\item Re-running this stage may overwrite previous metrics for the same \code{policy_id}.
\end{itemize}

\end{document}
